% gitver: кладёт git-версию в метаданные PDF и в \date (главный файл).
% В CI (dante-ev/latex-action) git падает с "dubious ownership" (репозиторий
% принадлежит хост-пользователю, а lualatex работает root'ом внутри
% контейнера), и gitver отвечает \PackageError --- строкой "!" в логе. В
% nonstopmode сборка продолжается, но latex-action считает любую строку "!"
% ошибкой шага. Поэтому в CI gitver не загружаем, а вместо него даём заглушку
% \gitVer. CI определяем по переменной окружения CI (её задаёт GitHub Actions
% в контейнере); читаем через Lua (\directlua), т.к. пакет env в образе
% dante-ev/texlive отсутствует. Обратный слэш внутри \directlua собираем через
% string.char(92) (см. комментарий про tikz externalization в главном файле).
\directlua{
  local bs = string.char(92)
  if os.getenv("CI") == nil then
    tex.sprint(bs .. "usepackage[noheader]{gitver}")
  else
    tex.sprint(bs .. "newcommand" .. bs .. "gitVer{CI build}")
  end
}
\usepackage{xurl} % Разрешить переносить URL на любой букве
\usepackage{caption}
\usepackage{subcaption}
\usepackage{morefloats}
\usepackage{float}

\NewDocumentCommand{\email}{m}{\href{mailto:#1}{#1}} % Кликабельный email

%% : с чуть более узки расстоянием по бокам.
%% Определено в pseudo, потому пока что закомментировано
 \NewDocumentCommand \rng { } {
 \nolinebreak
 \mathinner { : }
 \nolinebreak
 }

\newcommand{\bfgray}[1]{\cellcolor{lightgray}\textbf{#1}}
